\section{Giao diện quản lý}

Web Dashboard (\path{apps/dashboard}) là ứng dụng một trang viết bằng React 19
và TypeScript, build bằng Vite và định kiểu bằng Tailwind CSS. Mã nguồn được tổ
chức theo module chức năng; mỗi module gồm \texttt{pages}, \texttt{components},
\texttt{hooks}, \texttt{services} và \texttt{types}. Lớp \texttt{core/realtime} quản
lý kết nối WebSocket và đăng ký topic; \texttt{shared/streaming} xử lý luồng SSE
của Assistant. Khi triển khai, ứng dụng được phục vụ bởi Nginx, đồng thời Nginx
định tuyến \texttt{/api/v1/ws} (nâng cấp WebSocket, timeout 1 giờ) và
\texttt{/api/} tới Cloud, \texttt{/api/v1/chat} tới Assistant (tắt bộ đệm để
phát SSE), các đường dẫn còn lại trả về \texttt{index.html} theo cơ chế SPA. Trình
duyệt chỉ làm việc với một origin nên không phát sinh vấn đề CORS.

\begin{table}[H]
\centering
\caption{Các màn hình chính của Web Dashboard}
\label{tab:dashboard-pages}
\begin{tabularx}{\textwidth}{L{3cm}L{4.4cm}Y}
\toprule
\textbf{Module} & \textbf{Màn hình} & \textbf{Chức năng}\\
\midrule
overview & Overview, Live Monitor, Analytics & Tổng quan số cluster, node online/offline, runtime; theo dõi trực tiếp FPS, latency, tài nguyên\\
cluster & Clusters, Cluster Detail & Danh sách và chi tiết cluster cùng các node thành viên\\
agent & Edge Devices, Node Detail, Node Logs & Thông tin phần cứng, trạng thái Agent, runtime đang chạy, nhật ký sự kiện\\
agent & Cameras, Camera Groups, Snapshots & Quản lý nguồn camera cấp cho runtime và ảnh chụp kết quả\\
ai & Models, Class Config & Danh sách model và cấu hình lớp đối tượng của model\\
assistant & AI Assistant, Chat & Hội thoại với AI Assistant, hiển thị trạng thái gọi tool, biểu đồ và câu hỏi gợi ý\\
auth & Login & Đăng nhập\\
\bottomrule
\end{tabularx}
\end{table}

Màn hình Node Detail kết hợp dữ liệu REST (thông tin node, runtime) với luồng
WebSocket đăng ký \path{edge/<cluster>/<node>/telemetry/+} và
\path{edge/<cluster>/<node>/event/+}, nên biểu đồ và nhật ký cập nhật ngay khi Agent
báo cáo mà không cần tải lại trang. Màn hình AI Assistant hiển thị lần lượt các
bước ``đang phân tích'', ``đang gọi công cụ'', văn bản trả lời dần theo token, biểu
đồ do nút genchart sinh và các nút câu hỏi gợi ý.

Lớp \texttt{EdgeOpsContext} có chế độ mô phỏng (mock mode): khi bật,
\texttt{MockRealtimeFeed} sinh các bản tin cùng định dạng envelope với WebSocket
thật và đẩy vào cùng đường xử lý. Chế độ này cho phép phát triển và trình diễn
giao diện khi chưa có thiết bị Edge kết nối, mà không phải sửa mã các màn hình.
